\documentclass[10pt,a4paper]{article}
\usepackage[margin=22mm]{geometry}
\usepackage{amsmath,booktabs,siunitx,array,longtable}
\usepackage[T1]{fontenc}
\usepackage{lmodern}
\usepackage{hyperref}
\sisetup{group-separator={,}}
\setlength{\parskip}{4pt}
\setlength{\parindent}{0pt}
\title{S\&W framework refresh: kij refit on the S\&W 1992 equilibrium procedure\\[2pt]\large Coefficients, fit quality and comparison with the published set}
\author{Mark Burgoyne}
\date{11 November 2026 (helium section added 12 November 2026)}
\begin{document}
\maketitle

\section{Summary}

The correlations published in Burgoyne and Nielsen (2026), \emph{Fluid Phase Equilibria} 114824, were regressed with a binary gas-water solve that treats the vapour as pure, dry gas at total pressure. Søreide and Whitson (1992) fitted their correlations with one equilibrium in which the aqueous phase uses $k_{ij}^{\mathrm{AQ}}$ and the non-aqueous phase uses $k_{ij}^{\mathrm{NA}}$, with the vapour carrying its equilibrium water (their Eqs.~10a/10b). This note refits the aqueous-phase BIPs, re-checks the non-aqueous BIPs and refits the embedded salinity correction on that procedure, using the same data, exclusions, temperature window, functional forms and fitting routines as the paper. Only the equilibrium procedure changes. The published coefficients are unchanged and remain correct only inside the dry-vapour solve they were fitted with.

Findings, freshwater solubility on the paper's point set (Table~\ref{tab:fresh}):
\begin{itemize}
\item In a procedure-1 calculation the published $k_{ij}^{\mathrm{AQ}}$ underpredict solubility wherever the vapour holds appreciable water: bias $-12$ to $-29$\% at 150-200\,\textdegree C for CO$_2$, H$_2$S, CH$_4$ and H$_2$. The refit removes it ($-6$ to $+3$\%).
\item The refit shifts every pointwise $k_{ij}^{\mathrm{AQ}}$ downward (1,561 of 1,561 points), by less than 0.001 at or below 60\,\textdegree C and by 0.013-0.086 (gas mean) at 150-200\,\textdegree C, several times the pointwise scatter (MAE 0.006-0.035). Measurement noise cancels in this paired difference.
\item Scored on procedure 1, the S\&W 1992 correlations are closer to the refit than the paper reported: CO$_2$ 15.3\% vs 12.4\%, H$_2$S 8.0\% vs 6.8\% (excluding 18 suspect points, Section~\ref{sec:data}), CH$_4$ 3.9\% vs 3.9\%. The paper's S\&W column (23.0, 51.3, 5.8\%) scored S\&W's coefficients in the dry-vapour solve, which penalised them.
\item $k_{ij}^{\mathrm{NA}}$: no change justified for any gas under the paper's 1.1~pp rule. All published constants retained.
\item Helium (a post-publication extension, not in the paper): the procedure-1 refit moves $k_{ij}^{\mathrm{AQ}}$ by $-0.0002$ on average, because the helium correlation was already fitted on a gas-fugacity basis; freshwater MARE 6.3\% ($n = 130$), brine 12.8\% ($n = 48$). Section~\ref{sec:he}.
\item Brine: the refitted salt correction differs from the published one by 0.4-2.9\% in mean salt factor; brine MAREs are within the scatter of the published set once the freshwater bias is accounted for (Section~\ref{sec:brine}).
\end{itemize}

\section{Procedure and method}

Three ways of computing gas solubility with the two S\&W BIP sets appear in practice. They differ only in how much water the non-aqueous phase carries when the aqueous gas fugacity is equated to the vapour's:

\begin{center}
\begin{tabular}{p{38mm}p{62mm}p{52mm}}
\toprule
Procedure & Non-aqueous phase & Used by \\
\midrule
0, dry vapour & pure gas at total pressure & published paper (regression and MARE) \\
1, S\&W Eqs.~10a/10b & one equilibrium; aqueous phase $k^{\mathrm{AQ}}$, non-aqueous phase $k^{\mathrm{NA}}$ & S\&W 1992; this refit \\
2, dual flash & $k^{\mathrm{AQ}}$ in every phase for the solubility flash & Whitson \& Brulé (2000) shortcut for single-BIP EOS software \\
\bottomrule
\end{tabular}
\end{center}

Whitson and Brulé (2000, \S9.3) give procedure 2 for cases where only gas solubility or only water content is needed and note that predicting both requires an EOS implementation with phase-specific BIPs, which is procedure 1.

\textbf{EOS.} PR79, S\&W water alpha with NaCl molality (S\&W Eq.~9), applied to water in both phases (one pure-component set; only the BIPs are phase-specific). Component properties: S\&W Table~3 (H$_2$ and He from NIST), as in the paper.

\textbf{Solver.} Successive substitution on $K_i = \phi_i^{L}(x; k^{\mathrm{AQ}})/\phi_i^{V}(y; k^{\mathrm{NA}})$ with the closed-form binary Rachford-Rice, falling back to a nested bracketed solve. Verification against an independent implementation (Younus 2026, Table~4, seven binaries with S\&W 1992 BIPs): water contents for both procedures 1 and 2 agree to all printed digits, and the procedure-2/procedure-1 solubility ratios agree to $\pm0.01$ percentage points; fugacity residuals are below $2\times10^{-13}$ in $\ln f$.

\textbf{Regression.} Pointwise $k_{ij}^{\mathrm{AQ}}$ per freshwater point with $k_{ij}^{\mathrm{NA}}$ at its recommended value; correlations fitted by L1 (the paper's routine and forms); $k_{ij}^{\mathrm{NA}}$ re-checked with the paper's grid search and rows; embedded $\Delta k_{ij}$ fitted with the paper's routine to the same Sechenov targets (Dubessy et al.\ 2005 $-0.011$ for CO$_2$, Akinfiev 2016 $+0.019$ for H$_2$S, S\&W Eq.~8 otherwise) on the same $T$/$P$/$m$ grids. Because $k_{ij}^{\mathrm{NA}}$ did not change, no further $k^{\mathrm{AQ}}$/$k^{\mathrm{NA}}$ iteration was needed.

\section{Coefficients}

\subsection{Aqueous-phase BIP, freshwater}
Forms as published; $T$ in K, $T_r = T/T_c$.
\begin{align*}
\text{cubic:}\quad & k = c_0 + c_1 T + c_2 T^2 + c_3 T^3 \\
\text{exponential:}\quad & k = a/T + b\,\exp(c/T) + d \\
\text{linear:}\quad & k = a + bT \\
\text{rational:}\quad & k = (A + T_r)/(B + C\,T_r)
\end{align*}

\begin{center}\footnotesize\setlength{\tabcolsep}{3pt}
\begin{tabular}{llp{78mm}rrrl}
\toprule
Gas & Form & Coefficients & $n$ & MAE $k$ & max $|r|$ & $T$ fit (\textdegree C) \\
\midrule
CO$_2$ & cubic & $c_0{=}-1.550211$, $c_1{=}9.32272{\times}10^{-3}$, $c_2{=}-1.977702{\times}10^{-5}$, $c_3{=}1.527732{\times}10^{-8}$ & 606 & 0.0087 & 0.164 & 0-200 \\
H$_2$S & exp & $a{=}-82.78098$, $b{=}1.739491{\times}10^{6}$, $c{=}-8585.063$, $d{=}0.2518306$ & 405 & 0.0070 & 0.204 & 21-200 \\
CH$_4$ & rational & $A{=}-2.197678$, $B{=}0.7150154$, $C{=}0.894503$; $T_c{=}190.60$ & 115 & 0.0076 & 0.388 & 25-171 \\
N$_2$ & linear & $a{=}-1.621096$, $b{=}3.295273{\times}10^{-3}$ & 127 & 0.0104 & 0.034 & 25-200 \\
H$_2$ & rational & $A{=}-16.01983$, $B{=}0.7312363$, $C{=}0.681362$; $T_c{=}33.145$ & 154 & 0.0350 & 0.263 & 0-179 \\
C$_2$H$_6$ & rational & $A{=}-1.271348$, $B{=}-0.02927137$, $C{=}1.637183$; $T_c{=}305.40$ & 95 & 0.0063 & 0.057 & 2-200 \\
C$_3$H$_8$ & rational & $A{=}-1.156778$, $B{=}0.208946$, $C{=}1.810663$; $T_c{=}369.80$ & 59 & 0.0015 & 0.024 & 6-149 \\
He & rational & $A{=}-71.6155$, $B{=}-0.164407$, $C{=}0.5017371$; $T_c{=}5.1953$ & 130 & 0.0184 & 0.124 & 20-163 \\
\bottomrule
\end{tabular}
\end{center}

The H$_2$S exponential and the He and CH$_4$ rationals have strongly correlated parameters: individual coefficients moved a long way from the published values (H$_2$S $b$ from $1.349\times10^{3}$ to $1.739\times10^{6}$) while the curves moved by the amounts in Section~\ref{sec:shift}. Use all digits given. Heavier hydrocarbons (iC$_4$ and above) keep S\&W Eqs.~11-12, which S\&W fitted on procedure 1.

Check values, freshwater $k_{ij}^{\mathrm{AQ}}$ at 25 / 80 / 150 / 200\,\textdegree C: CO$_2$ $-0.12379$ / $-0.05152$ / $+0.01103$ / $+0.05158$; H$_2$S $-0.02582$ / $+0.01747$ / $+0.05889$ / $+0.09980$; CH$_4$ $-0.29959$ / $-0.14536$ / $+0.00830$ / $+0.09700$; N$_2$ $-0.63861$ / $-0.45737$ / $-0.22670$ / $-0.06194$; H$_2$ $-1.02393$ / $-0.67140$ / $-0.34499$ / $-0.16683$; C$_2$H$_6$ $-0.18807$ / $-0.06170$ / $+0.05101$ / $+0.11085$; C$_3$H$_8$ $-0.21005$ / $-0.10412$ / $-0.00549$ / $+0.04858$; He $-0.49694$ / $-0.10726$ / $+0.24159$ / $+0.42735$. Equilibrium check values (procedure 1, $x_{\mathrm{gas}}$ and $y_{\mathrm{H_2O}}$ at three states per gas) are in the accompanying \texttt{check\_values.csv}.

\subsection{Non-aqueous-phase BIP}
Unchanged from the paper: CO$_2$ 0.1896, H$_2$S 0.161, CH$_4$ 0.4850, N$_2$ 0.4778, H$_2$ 0.468, C$_2$H$_6$ 0.4920, C$_3$H$_8$ 0.5525, $n$C$_4$ 0.5091, C$_5$+ 0.500; He 0.468 (H$_2$ surrogate, no He water-content data). Table~\ref{tab:kna} gives the procedure-1 evidence.

\subsection{Embedded salinity correction}
$k_{ij}^{\mathrm{AQ}}(T,m) = k_{ij}^{\mathrm{AQ}}(T) + (a_0 + a_1T_r + a_2T_r^2)\,m\;[+\,(b_0 + b_1T_r)\,m^2$ for CO$_2$], $m$ = NaCl molality. The S\&W brine water alpha is applied with it. MAE and max are against the Sechenov targets on the fitting grid.

\begin{center}\footnotesize\setlength{\tabcolsep}{4pt}
\begin{tabular}{lrrrrrrrr}
\toprule
Gas & $T_c$ (K) & $a_0$ & $a_1$ & $a_2$ & $b_0$ & $b_1$ & MAE $x$ (\%) & max (\%) \\
\midrule
CO$_2$ & 304.20 & 0.0408481 & $-0.0789861$ & 0.0502757 & 0.00676611 & $-0.00735569$ & 2.92 & 15.8 \\
H$_2$S & 373.20 & 0.0474380 & $-0.0916139$ & 0.0498240 & & & 2.35 & 19.1 \\
CH$_4$ & 190.60 & 0.134391 & $-0.134258$ & 0.0408428 & & & 0.98 & 3.7 \\
N$_2$ & 126.10 & 0.219484 & $-0.148600$ & 0.0306167 & & & 1.04 & 4.2 \\
H$_2$ & 33.145 & 0.376712 & $-0.0647304$ & 0.00315541 & & & 0.66 & 2.4 \\
C$_2$H$_6$ & 305.40 & 0.0777696 & $-0.123360$ & 0.0625866 & & & 1.09 & 3.6 \\
C$_3$H$_8$ & 369.80 & 0.0580567 & $-0.111055$ & 0.0692266 & & & 1.40 & 3.3 \\
$n$C$_4$H$_{10}$ & 425.20 & 0.0456416 & $-0.100089$ & 0.0729267 & & & 1.75 & 3.4 \\
He & 5.1953 & 0.363020 & $-0.00963936$ & $7.09302{\times}10^{-5}$ & & & 0.56 & 1.3 \\
\bottomrule
\end{tabular}
\end{center}

C$_3$H$_8$ and $n$C$_4$H$_{10}$ are now fitted on their own freshwater BIPs (the paper carried them over from the MC-3 fit). The CO$_2$ and H$_2$S maxima sit at high temperature, at the 50~bar and 500~bar ends of the grid: on procedure 1 a pressure-independent $\Delta k_{ij}$ cannot follow the change in vapour water content with pressure as closely as it could in the dry-vapour solve (the published fit gave 1.56\% and 1.99\% MAE).

\section{Fit quality}

\subsection{Freshwater solubility}\label{sec:fresh}
Columns: \textbf{Pub-P0} published coefficients in the published dry-vapour solve (reproduces the paper exactly); \textbf{Pub-P1} published coefficients in a procedure-1 calculation; \textbf{SW-P1} S\&W 1992 aqueous and non-aqueous BIPs (Eq.~17 for H$_2$S) on procedure 1, as they were fitted; \textbf{New-P1} this refit. Evaluation set is the paper's (exclusions, 0-200\,\textdegree C; H$_2$ excludes Gillespie 1980 as in the paper). Five CO$_2$ points cannot be evaluated on procedure 1 (Section~\ref{sec:data}).

\begin{table}[h]\centering\small
\caption{Freshwater $x_{\mathrm{gas}}$ MARE (\%), with bias (\%) in brackets.}\label{tab:fresh}
\begin{tabular}{lrcccc}
\toprule
Gas & $n$ & Pub-P0 & Pub-P1 & SW-P1 & New-P1 \\
\midrule
CO$_2$ & 611 (606) & 16.4 ($+10.7$) & 14.9 ($-1.0$) & 15.3 ($-6.3$) & \textbf{12.4} ($+4.8$) \\
H$_2$S & 405 & 38.5 ($+30.4$) & 38.2 ($+13.9$) & 30.8 ($+20.6$) & 31.9 ($+25.7$) \\
H$_2$S, excl.\ 18 suspect & 387 & 13.8 & 13.8 & 8.0 & \textbf{6.8} \\
CH$_4$ & 115 & 4.8 ($-0.3$) & 4.7 ($-3.0$) & 3.9 ($-0.6$) & \textbf{3.9} ($-0.7$) \\
N$_2$ & 127 & 5.7 ($+1.6$) & 5.6 ($-0.4$) & 7.2 ($+3.7$) & \textbf{5.3} ($+1.2$) \\
H$_2$ & 142 & \textbf{10.2} ($+3.0$) & 12.0 ($-1.1$) & - & 11.3 ($+2.3$) \\
C$_2$H$_6$ & 95 & 12.9 ($+4.6$) & 9.2 ($-2.4$) & 11.3 ($-5.5$) & \textbf{8.4} ($+0.8$) \\
C$_3$H$_8$ & 59 & 5.0 ($+1.4$) & 4.8 ($-1.2$) & 8.9 ($+4.8$) & \textbf{3.4} ($+1.2$) \\
He & 130 & (9.7)$^{a}$ & 6.3 ($+0.1$) & - & 6.3 ($+0.1$) \\
\bottomrule
\end{tabular}\\[2pt]
\footnotesize $^{a}$The He correlation was fitted and published on a gas-fugacity basis (helium partial pressure), which is already close to procedure 1; its dry-vapour score is shown for completeness only.
\end{table}

The temperature trend is where the procedure matters. MARE (bias) at 150-200\,\textdegree C:

\begin{center}\small
\begin{tabular}{lrcccc}
\toprule
Gas & $n$ & Pub-P0 & Pub-P1 & SW-P1 & New-P1 \\
\midrule
CO$_2$ & 113 & 30.0 ($+23.3$) & 20.7 ($-20.4$) & 8.4 ($-6.0$) & 9.0 ($-2.4$) \\
H$_2$S & 103 & 31.9 ($+20.7$) & 28.8 ($-28.8$) & 8.1 ($+2.2$) & 6.8 ($+2.9$) \\
CH$_4$ & 21 & 12.7 ($-2.9$) & 12.8 ($-12.2$) & 9.3 ($-5.3$) & 9.2 ($-5.9$) \\
N$_2$ & 11 & 6.9 ($+3.5$) & 5.8 ($-5.8$) & 5.3 ($-5.3$) & 3.3 ($-0.8$) \\
H$_2$ & 20 & 9.1 ($+3.4$) & 17.2 ($-13.3$) & - & 11.6 ($-2.9$) \\
C$_2$H$_6$ & 26 & 14.1 ($+8.5$) & 6.9 ($-5.2$) & 6.0 ($-3.2$) & 5.8 ($+0.9$) \\
\bottomrule
\end{tabular}
\end{center}

At or below 60\,\textdegree C the published and refitted columns agree within 2~pp for every gas (H$_2$S aside, where the suspect points of Section~\ref{sec:data} dominate): the vapour is nearly dry and the procedure barely matters. H$_2$ is the one gas where the refit scores worse overall than the published solve (11.3\% vs 10.2\%); the difference is at 100-200\,\textdegree C, where the H$_2$ data are sparse (47 points) and the pointwise scatter is the largest of any gas (MAE 0.035).

\subsection{Size of the shift}\label{sec:shift}
Pointwise $k_{ij}^{\mathrm{AQ}}$ (procedure 1 minus published, same measured $x$), gas mean by temperature band:

\begin{center}\small
\begin{tabular}{lrrrr}
\toprule
Gas & $\le$60\,\textdegree C & 60-100 & 100-150 & 150-200 \\
\midrule
CO$_2$ & $-0.0008$ & $-0.0040$ & $-0.0114$ & $-0.0528$ \\
H$_2$S & $-0.0004$ & $-0.0021$ & $-0.0109$ & $-0.0483$ \\
CH$_4$ & $-0.0001$ & $-0.0006$ & $-0.0043$ & $-0.0185$ \\
N$_2$ & $-0.0003$ & $-0.0042$ & $-0.0061$ & $-0.0266$ \\
H$_2$ & $-0.0003$ & $-0.0083$ & $-0.0330$ & $-0.0863$ \\
C$_2$H$_6$ & $-0.0007$ & $-0.0006$ & $-0.0061$ & $-0.0126$ \\
C$_3$H$_8$ & $-0.0001$ & $-0.0008$ & $-0.0040$ & - \\
\bottomrule
\end{tabular}
\end{center}

\subsection{Water content}
\begin{table}[h]\centering\small
\caption{$y_{\mathrm{H_2O}}$ MARE (\%) on procedure 1 with the refitted $k_{ij}^{\mathrm{AQ}}$; rows as the paper's $k_{ij}^{\mathrm{NA}}$ analysis.}\label{tab:kna}
\begin{tabular}{lrrrrrr}
\toprule
Gas & $n$ & $k^{\mathrm{NA}}$ & MARE & S\&W value & optimum $k^{\mathrm{NA}}$ (MARE) & paper (dry solve, MC-3 alpha) \\
\midrule
CO$_2$ & 208 & 0.1896 & 31.7 & 31.7 & 0.190 (31.7) & 35.1 \\
H$_2$S & 162 & 0.161 & 6.7 & 7.1 (Eq.~17) & 0.147 (6.4) & 8.6 \\
CH$_4$ & 93 & 0.4850 & 7.1 & 7.1 & 0.495 (6.7) & 10.2 \\
N$_2$ & 31 & 0.4778 & 8.1 & 8.1 & 0.537 (7.5) & 7.2 \\
H$_2$ & 14 & 0.468 & 4.2 & - & 0.499 (4.1) & 4.5 \\
C$_2$H$_6$ & 58 & 0.4920 & 19.0 & 19.0 & 0.476 (18.1) & 21.1 \\
C$_3$H$_8$ & 11 & 0.5525 & 5.5 & 5.5 & 0.557 (5.4) & 4.1 \\
\bottomrule
\end{tabular}
\end{table}
The largest re-optimisation gain is 0.9~pp (C$_2$H$_6$), below the paper's 1.1~pp threshold, so every constant stays. On procedure 1 the paper's H$_2$S change (0.161 replacing S\&W Eq.~17) gains 0.44~pp against 3.1~pp in the dry solve; it is retained for continuity, but by the paper's own rule it would no longer be justified.

\subsection{Brine}\label{sec:brine}
\begin{center}\small
\begin{tabular}{lrcccc}
\toprule
Gas & $n$ & Pub-P0 & Pub-P1 & SW-P1 & New-P1 \\
\midrule
CO$_2$ & 109 & 17.3 ($+10.0$) & 18.0 ($+7.3$) & 18.4 ($-15.2$) & 17.1 ($+11.6$) \\
H$_2$S & 100 & 15.4 ($+8.1$) & 7.6 ($-2.7$) & 64.9 ($-64.9$) & 11.0 ($+4.3$) \\
CH$_4$ & 32 & 5.6 ($+2.4$) & 5.1 ($+1.7$) & 11.9 ($-9.8$) & 6.5 ($+3.3$) \\
N$_2$ & 45 & 4.6 ($-2.8$) & 4.8 ($-3.3$) & 4.3 ($-1.3$) & 4.2 ($-1.5$) \\
H$_2$ & 47 & 4.1 ($-0.3$) & 3.9 ($-0.8$) & - & 5.8 ($-2.5$) \\
He & 48 & 13.0 ($+3.4$) & 12.8 ($+3.2$) & - & 12.8 ($+3.2$) \\
\bottomrule
\end{tabular}
\end{center}

The published coefficients score better on H$_2$S, CH$_4$ and H$_2$ brine in a procedure-1 calculation than the refit does. Splitting each prediction into its freshwater baseline and its salt factor shows why: the mean salt factors of the two sets agree within 0.4-2.9\% per gas, and the difference comes from the freshwater baseline. For Suleimenov (1994) H$_2$S at 155\,\textdegree C the published baseline is 28\% low on procedure 1, which offsets a salt correction that is too weak at that temperature (brine error $-9.7$\%); with the corrected baseline the refitted salt correction gives $+17.1$\%. The refit exposes a high-temperature weakness in the H$_2$S Sechenov model rather than introducing an error.

Sensitivity: defining the Sechenov target as an activity coefficient at equal gas fugacity, instead of the paper's ratio at equal $T$ and total $P$ (which is also how the measured $k_s$ were extracted), changes every brine MARE above by at most 0.3~pp. The total-pressure definition is used.

\section{Helium}\label{sec:he}

Helium was added to the framework after publication (October 2026) and is not part of the peer-reviewed paper. It is documented here in full because the procedure-1 basis is the first one on which its coefficients are recommended alongside the other gases.

\textbf{Properties.} NIST (Ortiz-Vega et al.\ 2019 equation of state): $T_c = 5.1953$~K, $P_c = 0.228323$~MPa, $\omega = -0.38354$, $T_b = 4.2238$~K.

\textbf{Data.} 130 freshwater points from five sources, 20-163\,\textdegree C and 1-1,000~bar, in \texttt{shared/data/he\_solubility.csv} with DOIs (fit and validate flags as in the October fit): Abrosimov and Lebedeva (2013, evaluated values; 105 points, 50-80\,\textdegree C, 1-1,000~bar), Gardiner and Smith (1972; 12, 20-50\,\textdegree C), Gerth (1983; 8, 25\,\textdegree C), Potter and Clynne (1978; 2, 104-149\,\textdegree C) and Pray et al.\ (1952; 3, 163\,\textdegree C). The October fit targeted pointwise $k_{ij}$ at the PR fugacity of pure helium at the helium partial pressure, which already approximates procedure 1; the vapour water fraction on procedure 1 ranges from $1.5\times10^{-4}$ to 0.50 across these points.

\textbf{Aqueous-phase BIP.} Rational form, coefficients in Section~3.1 ($A = -71.6155$, $B = -0.164407$, $C = 0.5017371$). The individual coefficients differ from the October values ($-71.6152$, $-0.394962$, $0.505807$), but the curve moves by $-0.0002$ on average ($-0.0002$ to $+0.0008$) over the data: the rational parameters are strongly correlated.

\begin{center}\small
\begin{tabular}{lrccrrr}
\toprule
Source & $n$ & $T$ (\textdegree C) & $P$ (bar) & Pub-P0 & New-P1 (bias) & max \\
\midrule
Abrosimov \& Lebedeva 2013 & 105 & 50-80 & 1-1,001 & 9.5 & 6.8 ($+0.2$) & 12.6 \\
Gardiner \& Smith 1972 & 12 & 20-50 & 101-608 & 2.4 & 2.4 ($+0.7$) & 6.2 \\
Gerth 1983 & 8 & 25 & 51-405 & 2.8 & 2.8 ($-0.3$) & 5.3 \\
Potter \& Clynne 1978 & 2 & 104-149 & 11-13 & 39.0 & 2.7 ($+2.7$) & 2.9 \\
Pray et al.\ 1952 & 3 & 163 & 14-27 & 46.2 & 15.9 ($-6.2$) & 27.0 \\
\bottomrule
\end{tabular}
\end{center}

MARE in \%. The large Pub-P0 errors at low pressure and high temperature are the dry-vapour solve, not the correlation: at 104-163\,\textdegree C and 11-27~bar the vapour is 10-50\% water. On Abrosimov and Lebedeva the residuals carry a pressure trend (overprediction at 1-10~MPa, underprediction at 70-100~MPa), the structured bias of a $k_{ij}(T)$-only form. Not validated above 163\,\textdegree C.

\textbf{Non-aqueous-phase BIP.} 0.468, the H$_2$ value used as a surrogate. No helium water-content data were found, so water content of helium-rich gas is unvalidated.

\textbf{Embedded salinity correction.} Coefficients in Section~3.3, fitted to S\&W Eq.~8 Sechenov targets ($T_b = 4.2238$~K) on the H$_2$ grid restricted to 25-100\,\textdegree C. Eq.~8 matches measured helium $k_s$ in NaCl to 0.0067~kg/mol mean absolute deviation ($n = 48$). Brine MARE is 12.8\% ($n = 48$, bias $+3.2$\%), dominated by Gardiner and Smith's 100\,\textdegree C polynomial points ($+20$ to $+43$\%), which Potter and Clynne show are wrong at that temperature.

\section{Data notes}\label{sec:data}
\begin{itemize}
\item CO$_2$, Wang (2021), 5 points at 160-200\,\textdegree C and 5-10~bar: the stated pressure is at or below the water vapour pressure (6.2~bar at 160\,\textdegree C, 15.5~bar at 200\,\textdegree C), so it must be CO$_2$ partial pressure. No aqueous phase exists on procedure 1 at that total pressure; excluded from the fit and the procedure-1 MAREs.
\item H$_2$S, Selleck (1952) at 37.8\,\textdegree C, 6 points: tabulated $x_{\mathrm{H_2S}}$ falls from 0.0037 to 0.0011 as pressure rises from 7 to 27~bar, which no solubility isotherm does; every model is off by about 1,200\%. Suleimenov (1994) at or below 60\,\textdegree C, 12 points, about 250\% for every model. These 18 points are in the paper's set and set its H$_2$S MARE (38.5\%; 13.8\% without them). Flagged for source check, not removed here.
\item Nine H$_2$S points near the water-saturated mixture critical locus (Selleck 1952, Koschel 2013; 104-138\,\textdegree C, 155-300~bar) need an edge-refined root search on procedure 1; all converge.
\end{itemize}

\section{Limitations}
\begin{itemize}
\item Binary data only, as in the paper. Procedure 1 makes the coefficients transferable to multicomponent and oil systems in principle; that is not tested against data here.
\item The forms and the temperature limits are the paper's; the correlations should not be extrapolated beyond the fitted ranges listed above (200\,\textdegree C at most).
\item Brine targets rely on the paper's Sechenov models; the H$_2$S model is weak at 155\,\textdegree C (Section~\ref{sec:brine}).
\end{itemize}

\section*{References}
\small
Burgoyne, M., Nielsen, M.H., 2026. \emph{Fluid Phase Equilibria} 114824, doi:10.1016/j.fluid.2026.114824.\\
Abrosimov, V.K., Lebedeva, E.Yu., 2013. \emph{Russ.\ J.\ Inorg.\ Chem.} 58, 808-812, doi:10.1134/S0036023613070024.\\
Gardiner, G.E., Smith, N.O., 1972. \emph{J.\ Phys.\ Chem.} 76, 1195-1202, doi:10.1021/j100652a019.\\
Gerth, W.A., 1983. \emph{J.\ Solution Chem.} 12, 655-669, doi:10.1007/BF00648669.\\
Ortiz-Vega, D.O., Hall, K.R., Holste, J.C., Arp, V.D., Harvey, A.H., Lemmon, E.W., 2019. Helium equation of state (NIST).\\
Potter, R.W., Clynne, M.A., 1978. \emph{J.\ Solution Chem.} 7, 837-844, doi:10.1007/BF00650811.\\
Pray, H.A., Schweickert, C.E., Minnich, B.H., 1952. \emph{Ind.\ Eng.\ Chem.} 44, 1146-1151, doi:10.1021/ie50509a058.\\
Søreide, I., Whitson, C.H., 1992. \emph{Fluid Phase Equilibria} 77, 217-240 (with erratum).\\
Whitson, C.H., Brulé, M.R., 2000. Phase Behavior. SPE Monograph 20, Ch.~9.\\
Younus, B., 2026. Water PVT with the Søreide-Whitson model. Whitson internal study (November 2026).

\end{document}
